% from ega3-1-fr.tex
\oldpage[III]{145}

\subsection{Un critère d'amplitude}

\begin{theorem}[4.7.1]
\label{III.4.7.1-fr}
Soient $Y$ un préschéma localement noethérien, $f:X\to Y$ un morphisme
propre, $\mathcal L$ un $\mathcal O_X$-Module inversible, $y$ un point de
$Y$, $X_y=X\otimes_Yk(y)=f^{-1}(y)$, $g$ la projection de $X_y$ dans $X$.
Si $\mathcal L_y=g^*(\mathcal L)=\mathcal L\otimes_{\mathcal O_Y}k(y)$ est
ample sur $X_y$, il existe un voisinage ouvert $U$ de $y$ dans $Y$ tel que
$\mathcal L|f^{-1}(U)$ soit ample pour la restriction de $f$ à $f^{-1}(U)$.
\end{theorem}

\begin{proof}
\noindent\textit{I)} Posons
\[
Y'=\operatorname{Spec}(\mathcal O_y),\qquad
X'=X\times_YY',\qquad
\mathcal L'=\mathcal L\otimes_{\mathcal O_Y}\mathcal O_y~;
\]
nous allons d'abord prouver que $\mathcal L'$ est ample pour
$f'=f_{(Y')}$. On a le diagramme commutatif
\[
\xymatrix{
X\ar[d]_f & X'\ar[l]\ar[d]_{f'} & X_y\ar[l]\ar[d]_{f_y}\\
Y & Y'\ar[l] & \operatorname{Spec}(k(y))\ar[l]
}
\]
Comme $f'$ est propre (II, 5.4.2, (iii)) et $\mathcal O_y$ noethérien, on
voit qu'on peut se borner au cas où
$Y=Y'=\operatorname{Spec}(\mathcal O_y)$, donc $X=X'$, supposer que
$\mathcal L_y$ est ample pour $f_y$ et prouver que $\mathcal L$ est ample
pour $f$ (II, 4.6.6). Nous allons appliquer le critère
(\hyperref[III.2.6.1-fr]{2.6.1, c)}) et montrer en fait que pour tout
$\mathcal O_X$-Module cohérent $\mathcal F$, il existe un entier $N$ tel que
\[
H^1(X,\mathcal F(n))=0\quad\text{pour tout }n\geq N,
\qquad
\mathcal F(n)=\mathcal F\otimes\mathcal L^{\otimes n}.
\]
Notons que $y$ est un point fermé de $Y$ correspondant à l'idéal maximal
$\mathfrak m$ de $\mathcal O_y$~; $X_y$ est donc un sous-préschéma fermé de
$X$ défini par l'Idéal cohérent
$\mathcal J=f^*(\widetilde{\mathfrak m})\mathcal O_X=
\mathfrak m\mathcal O_X$ de $\mathcal O_X$ (I, 4.4.5), et $g$ l'injection
canonique. Considérons alors la $k(y)$-algèbre graduée
\[
S=\operatorname{gr}(\mathcal O_y)
=\bigoplus_{j\geq0}\mathfrak m^j/\mathfrak m^{j+1},
\]
qui est de type fini puisque $\mathcal O_y$ est noethérien~; la
$\mathcal O_X$-Algèbre $\mathcal S=f^*(\widetilde S)$ est donc
quasi-cohérente et de type fini, et elle est évidemment annulée par
$\mathcal J$, donc si on pose $\mathcal S_y=g^*(\mathcal S)$,
$\mathcal S_y$ est une $\mathcal O_{X_y}$-Algèbre quasi-cohérente de type
fini, et $\mathcal S=g_*(\mathcal S_y)$. Posons d'autre part
\[
\mathcal M_j=\mathfrak m^j\mathcal F/\mathfrak m^{j+1}\mathcal F,
\qquad
\mathcal M=\bigoplus_{j\geq0}\mathcal M_j
=\operatorname{gr}(\mathcal F)~;
\]
comme $\mathcal F$ est cohérent, $\mathcal M$ est un $\mathcal S$-Module
quasi-cohérent de type fini (0, 10.1.1) qui est aussi annulé par
$\mathcal J$, de sorte que si l'on pose
\[
\mathcal M'_j=g^*(\mathcal M_j),\qquad
\mathcal M'=g^*(\mathcal M)=\bigoplus_{j\geq0}\mathcal M'_j,
\]
$\mathcal M'$ est un $\mathcal S_y$-Module gradué quasi-cohérent de type
fini tel que $\mathcal M=g_*(\mathcal M')$. En outre, si on pose
$\mathcal M'_j(n)=\mathcal M'_j\otimes\mathcal L_y^{\otimes n}$, on a
$\mathcal M'_j(n)=g^*(\mathcal M_j(n))$. Cela étant, $f_y$ est propre
(II, 5.4.2, (iii)) et $\mathcal L_y$ est ample, donc $f_y$ est projectif
(II, 5.5.4 et 4.6.11), et on peut appliquer à
$\operatorname{Spec}(k(y))$, $f_y$, $\mathcal S_y$, $\mathcal L_y$ et
$\mathcal M'$ le théorème (\hyperref[III.2.4.1-fr]{2.4.1, (ii)})~: il
existe un entier $N$ tel que pour $n\geq N$, on ait
$H^q(X_y,\mathcal M'_j(n))=0$ pour tout $q>0$ et tout $j$~; par suite, on a
aussi $H^q(X,\mathcal M_j(n))=0$ pour tout $q>0$ et tout $j$
(G, II, 4.9.1). Posons alors
\[
\mathcal F(n)_j=\mathcal F(n)/\mathfrak m^{j+1}\mathcal F(n),
\]
de sorte que $\mathcal F(n)_{j-1}=\mathcal F(n)_j/\mathcal M_j(n)$ pour
$j\geq1$ et $\mathcal F(n)_0=\mathcal M_0(n)$. On a
$H^1(X,\mathcal F(n)_0)=0$, et, par la suite exacte de cohomologie,
$H^1(X,\mathcal F(n)_j)=H^1(X,\mathcal F(n)_{j-1})$ pour tout $j\geq1$,
donc $H^1(X,\mathcal F(n)_j)=0$ pour tout $j\geq0$. On conclut donc de
(\hyperref[III.4.2.1-fr]{4.2.1}) que $H^1(X,\mathcal F(n))=0$, ce qui
achève de prouver notre assertion.

